\section{Introdu\c{c}\~ao}

O rastreamento do c\^ancer do colo do \'utero depende de programas capazes de
identificar altera\c{c}\~oes precursoras antes da progress\~ao da doen\c{c}a. Em
2020, foram estimados aproximadamente 604 mil novos casos e 342 mil \'obitos no
mundo, com maior carga onde vacina\c{c}\~ao, rastreamento e tratamento s\~ao menos
acess\'iveis \cite{singh2023global}. Na citologia convencional, a automa\c{c}\~ao
tem sido estudada em classifica\c{c}\~ao de recortes, segmenta\c{c}\~ao,
detec\c{c}\~ao de inst\^ancias e an\'alise de campos ou l\^aminas
\cite{jiang2023systematic}.

Essas tarefas dependem do tipo de anota\c{c}\~ao dispon\'ivel. Um ponto informa
onde uma inst\^ancia est\'a, mas n\~ao sua extens\~ao; uma caixa acrescenta uma
regi\~ao aproximada; uma m\'ascara descreve fronteiras. Portanto, transformar
pontos em caixas n\~ao \'e apenas uma convers\~ao de formato: \'e uma decis\~ao
sobre qual regi\~ao o detector deve aprender. Em campos com muitas c\'elulas, essa
decis\~ao tamb\'em afeta sobreposi\c{c}\~ao entre alvos, localiza\c{c}\~ao e contagem
\cite{chen2023p2bnet,chen2021sparsepoints,xia2023colorization}.

A CRIC Cervix disponibiliza 400 imagens convencionais de Papanicolau e 11.534
c\'elulas classificadas, cada uma associada a uma coordenada nuclear
\cite{rezende2021cric}. A base permite estudar detec\c{c}\~ao sem recortar
previamente as c\'elulas, mas n\~ao fornece bounding boxes anat\^omicas. Assim, o
uso de detectores de objetos exige uma regra expl\'icita para converter cada
ponto em alvo de treinamento.

Este artigo formula o problema como detec\c{c}\~ao e contagem supervisionadas por
pontos. N\~ao se prop\~oe uma nova arquitetura nem se usa a categoria Bethesda
como sa\'ida: todas as coordenadas s\~ao tratadas como inst\^ancias de uma classe
\'unica. O objetivo \'e avaliar se pseudo-caixas centradas nos pontos nucleares
permitem recuperar as inst\^ancias anotadas e preservar a contagem por imagem.

Duas perguntas orientam o experimento. \textbf{RQ1:} como o tamanho da
pseudo-caixa afeta F1, AP50 e MAE de contagem, e qual tamanho deve ser levado ao
protocolo final? \textbf{RQ2:} com o tamanho selecionado, qual desempenho \'e
obtido em valida\c{c}\~ao cruzada 5-fold por imagem? A contribui\c{c}\~ao principal \'e separar a etapa
explorat\'oria de escolha do alvo da valida\c{c}\~ao final, tornando audit\'avel uma
decis\~ao que costuma ficar escondida na prepara\c{c}\~ao dos dados.
